\section{The three-sector bridge fibre}\label{sec:bridge-fibre}

Cut every bridge of the labelled reduced tree.  Its component-incidence graph
is a tree \(T\).  For a component \(v\), let \(P_v\) be its normalized Fourier
boundary tensor, including one character \(h_e\) for every incident bridge,
with \(P_v(0)=1\).  For \(e=uv\), write

\[
 k_e(0)=1,\qquad k_e(C)=c_e,\quad k_e(G)=g_e,\quad k_e(T)=t_e.
\]

Character conservation gives the global contraction

\begin{equation}\label{eq:bridge-contraction}
 Q(\mathbf h)=
 \prod_{v\in V(T)}P_v(\mathbf h|_v)
 \prod_{e\in E(T)}k_e(h_e).
\end{equation}

\begin{theorem}[Complete positive three-sector bridge fibre]
\label{thm:bridge-fibre}
On the positive normalized K3P component-tensor locus, the complete fibre of
\eqref{eq:bridge-contraction} is exactly

\begin{align}
 P_v(\mathbf h)&\longmapsto P_v(\mathbf h)
 \prod_{e\ni v}\prod_{s\in\{C,G,T\}}
 (a^s_{v,e})^{\1[h_e=s]},\label{eq:bridge-action-component}\\
 k_e(s)&\longmapsto
 \frac{k_e(s)}{a^s_{u,e}a^s_{v,e}}
 \qquad(s=C,G,T),\label{eq:bridge-action-edge}
\end{align}

with positive incidence scales and \(a^0_{v,e}=1\).  There is no additional
positive continuous or discrete gauge, the action is free on every retained
component, and the bridge tree has no holonomy.  On physical networks the
fibre is the intersection of this orbit with all local and edge inequalities.
\end{theorem}

\begin{proof}
Cut one bridge and fix \(s\in\{C,G,T\}\).  Its Fourier block is a positive
rank-one matrix, whose factorization is unique up to one positive scalar.
All-zero normalization fixes the zero-sector scalar.  Unlike K2P, the three
nonzero characters are observably labelled and independent; no equality or
permutation is imposed.  Peeling leaves of \(T\) sector by sector assigns one
scale to every incidence and gives
\eqref{eq:bridge-action-component}--\eqref{eq:bridge-action-edge}.  Direct
cancellation proves the converse, while the tree structure leaves no cycle
on which a holonomy could circulate.

Call a retained component \emph{marked} if it carries a retained physical
leaf boundary block.  For each incident bridge and nonzero sector, that block
supplies the compensating character in a conservation-supported
one-character anchor.  Call the component \emph{unmarked} otherwise.  Every
unmarked retained component in the strong class has bridge degree at least
three: the only possible degree-two theta stabilizer is the two-boundary
\(K_4-e\) factor, which has no tree-child rooting.

For freeness, a marked component has one-character anchors with identity
exponent matrix.  For an unmarked degree-\(d\) component, \(d\ge3\), use pair
anchors

\[
 (1,2),(1,3),(2,3),(1,4),\ldots,(1,d).
\]

In one sector the leading \(3\)-by-\(3\) exponent block has determinant
\(-2\), and every later row introduces one new coordinate.  Its positive
inverse is

\begin{equation}\label{eq:bridge-normalizer}
 a_1=\sqrt{\frac{r_{12}r_{13}}{r_{23}}},\quad
 a_2=\sqrt{\frac{r_{12}r_{23}}{r_{13}}},\quad
 a_3=\sqrt{\frac{r_{13}r_{23}}{r_{12}}},\quad
 a_k=\frac{r_{1k}}{a_1}.
\end{equation}

Three block-diagonal copies have rank \(3d\) and leading determinant
\((-2)^3=-8\).  Positivity makes the slice analytic.  The excluded unmarked
degree-two stabilizer cannot occur in a retained strong-class factor: the
two-boundary \(K_4-e\) theta has no tree-child rooting.  Finally, a
distinct-spectrum anchor such as
\((q_{CC},q_{GG},q_{TT})=(1/4,1/9,1/16)\) excludes a hidden permutation of
the fixed state labels.
\end{proof}

\begin{lemma}[Physical local product chart]\label{lem:physical-product}
Near every regular physical point, the quotient by the action in
\cref{thm:bridge-fibre} is a physical local product chart, simultaneously on
all bridges, in both \(\Dplus\) and \(\Dct\).
\end{lemma}

\begin{proof}
For one bridge spectrum \(y\), choose two isotropic endpoint factors near the
identity and a compensating residual factor, using
\cref{lem:subdivision}.  All six endpoint sector coordinates may vary
independently in a sufficiently small neighborhood while the residual stays
strict.  Do this for finitely many bridges and intersect the neighborhoods.
The analytic normalizers extract local projective component tensors and
effective bridge coordinates from the global tensor, while contraction is
their local inverse.  The same argument uses \(B_{\mathrm{CT}}\) on the
continuous-time cone.
\end{proof}

For a complete ported factor \(H\), quotient its normalized tensor by one
positive \(C,G,T\) scale per port and write \(\overline\Phi_H\) for a local
analytic slice map.  Local directed containment is defined exactly as in
\cref{def:relations}, using a source-open regular germ and a physical target
section.  Different positive anchor charts give equivalent definitions by
\cref{thm:bridge-fibre}.
